\section{Solutions for Chapter~\ref{ch:identification}}

\paragraph{Exercise~\ref{ex:identification-product}.}
\phantomsection
\label{sol:ex:identification-product}
Independent stress tensors make the dimension of a product
field the sum of its factor dimensions. The mixed three-point
function factors into two normalized two-point functions,
giving $\lambda=1$ in the displayed product convention.
Moreover $c_\psi=c_\phi c_\chi$, hence $I_\lambda=1$ whenever
its denominator cubics are nonzero.
Under independent rephasings, $\lambda$ itself can acquire a
phase, while $|\lambda|$, all dimensions, and $I_\lambda$
remain unchanged. Thus the invariant comparison is stronger
than comparing a printed phase of one matrix element.



\paragraph{Exercise~\ref{ex:identification-Narain}.}
\phantomsection
\label{sol:ex:identification-Narain}
On a totally isotropic subspace, projection to the positive
two-dimensional factor has zero kernel: a vector with positive
part zero must also have negative part zero by nullness. The
projection is injective, so the subspace has dimension at most
two. For light scalar charges the integral pairing has absolute
value strictly below one and must vanish. At a non-strict
endpoint the value $\pm1$ is not excluded, so the planar
reduction requires a separate argument.



\paragraph{Exercise~\ref{ex:identification-triad}.}
\phantomsection
\label{sol:ex:identification-triad}
The only nonzero charge-neutral ordered triples are the six
permutations of $(p_1,p_2,p_3)$. Each coefficient product is
$1/(\sqrt3)^3$, giving $2/\sqrt3$. All three external dimensions
are equal, exactly the case excluded by the unequal-shell
hypothesis. Arbitrary degeneracy within unequal shells was
already allowed by that theorem; it is equality of the three
radii that changes the geometry.



\paragraph{Exercise~\ref{ex:identification-crossing}.}
\phantomsection
\label{sol:ex:identification-crossing}
Write $q(z)=(1-z)^{2\delta}$ and $P(z)=q(z)g_D(z)$.
Then $F_D(z)=P(z)-P(1-z)$, so the first and third derivatives
at $1/2$ are respectively $2P'$ and $2P'''$. At that point,
\[
 \frac{q'}q=-4\delta,\qquad
 \frac{q''}q=8\delta(2\delta-1),\qquad
 \frac{q'''}q=-32\delta(2\delta-1)(\delta-1).
\]
Insert $g_D''=2g_D'+8Cg_D$ and
$g_D'''=8(C+1)g_D'$, where $C=D(D-1)$, into the product
rule. With the chapter's $A=1-6\delta+6\delta^2$ and
$B=1-3\delta+2\delta^2$, this gives
\[
 \frac{P'''}q=8(C+A)g_D'-32\delta(3C+B)g_D,
 \qquad \frac{P'}q=g_D'-4\delta g_D.
\]
Substitution into the defined functional yields
\[
 \frac{\mathcal L(F_D)}{g_D}
 =(C+A-\kappa)r_D-4\delta(3C+B-\kappa),
 \qquad r_D=g_D'/g_D.
\]
For the identity $D=0$, use $g_0=1$ directly rather than the
positive-dimension integral representation.

The arithmetic in the certificate is exact:
\[
 \frac54\frac9{10}(.083)=.093375,\qquad
 \frac{169}{100}\frac{37}{40}(.246)=.3845595,
\]
and hence $.303+.093375-.3845595=.0118155>0$.
The identity, charged, and epsilon contributions have a positive
total under the stated lower bounds. Every additional exchange
outside $(.16,1.75)$ has nonnegative coefficient and positive
functional value, so crossing requires another coupled exchange
inside that interval.

Application to the chain requires a unitary local continuum
CFT with the specified normalized Hermitian scalar and a positive,
convergent block expansion. The proposed dimensions and OPE
magnitudes must be limiting continuum values, or controlled
intervals over which the strict inequality survives. Finite
eigenvector overlaps alone do not establish those premises.
The additional interpretation as a new two-dimensional scalar
uses integer locality, even spin in the identical-scalar OPE,
and the descendant analysis in the chapter. It does not require
assuming a particular central charge or a named Haagerup CFT.
It also does not exclude the continuum when the required extra
scalar is present.



\paragraph{Exercise~\ref{ex:identification-orbifold}.}
\phantomsection
\label{sol:ex:identification-orbifold}
The oscillator formula gives
$h_L=11/288$ and $h_R=35/288$. Their sum is $23/144$,
but that pair is not scalar: its spin is $-1/12$.
Nonnegative excitations and zero modes imply that a scalar
must have both weights at least $35/288$, hence
$\Delta\geq35/144$. Level matching or projection can
increase this threshold; the bound need not be attained.

If $a^3=1$ in the abelianization and $\chi^2=1$, then
$\chi(a)^3=\chi(a)^2=1$, so $\chi(a)=1$. This applies
to each determinant character independently. For the fixed
Hermitian axis, the unit normalization
$X=(e^{i\theta}\phi+e^{-i\theta}\phi^\dagger)/\sqrt2$
gives $|C_{XXX}|=|c_\phi|/\sqrt2$ at a cubic extremum.
The isometric lift contributes no group-order factor.
Thus $|c_\phi|/\sqrt2\leq\sqrt{2/3}$ gives the stated
$2/\sqrt3$, with no further rescaling.



\paragraph{Exercise~\ref{ex:identification-Potts}.}
\phantomsection
\label{sol:ex:identification-Potts}
Choose the common Potts field phase so that $c_\sigma$ is real
and positive. The triples $w_i^3$ then contribute $c_\sigma^2a_i^3$:
two factors contribute conjugate order cubics, whose product is
$|c_\sigma|^2=c_\sigma^2$, and the third contributes the identity.
The six permutations of $w_1w_2w_3$ pair each
Potts order field with its adjoint and contribute
$6a_1a_2a_3$. Every other triple violates an independent
ultraviolet charge. In $C(\sigma_1,\chi,\sigma_2^\dagger)$
only $w_1$ survives, giving $\lambda=a_1$ and hence
$|\lambda|\leq1$. A large cubic can arise from several
charge-conserving permutations without increasing this one
mixed matrix element.



\paragraph{Exercise~\ref{ex:identification-norm}.}
\phantomsection
\label{sol:ex:identification-norm}
The conjugate of $\varphi$ is $(1-\sqrt5)/2$. Therefore
\[
 N(\varphi+2)=\frac{(5+\sqrt5)(5-\sqrt5)}4=5,
 \qquad
 N\!\left(\frac{64(\varphi+2)^3}{3}\right)
 =\frac{64^2 5^3}{9}.
\]
An algebraic integer has an integer field norm, so the second
number proves that the proposed quotient is not an algebraic
integer. Fusion-subcategory divisibility then excludes an
order-three pointed subcategory from the ambient category of
defects preserving the original retained algebra, and hence
from its smaller maximal-chiral-trivial symmetry kernel.
It does not control a symmetry that permutes stress tensors
and acts nontrivially on the chiral algebra. The separate
three-cycle argument, together with the maximal graph pairing,
is what bounds those chiral images and prevents independent
left-right cycles.



\paragraph{Exercise~\ref{ex:identification-residues}.}
\phantomsection
\label{sol:ex:identification-residues}
Reduce the sums modulo one, keep values strictly between zero
and $1/2$, and then double them. Adding zero to the Fibonacci
residues retains only $1/5$, yielding $\Delta=2/5$. Adding
$1/2$ retains $1/10$ and $3/10$, yielding $1/5$ and $3/5$.
Adding $5/16$ retains $5/16$ and $9/80$, yielding $5/8$ and
$9/40$. This gives the chapter's five-element list.

For a scalar below dimension one, positive-energy integer
grading leaves no freedom to add a nonzero integer to its
chiral weight while remaining below $1/2$. Therefore twist
residues impose this necessary condition. The full pairing
can remove possible entries, and the category does not fix
their actual multiplicities or lowest characters. An
intermediate smaller algebra has a different representation
category, so this residue list cannot be imposed on it by
performing a further mathematical condensation absent from
the physical theory.

